\section{A collision-entropy bound in every residue class}
\label{sec:residue-bound}
Let $d\ge2$ unbiased pairwise-independent sign variables have joint
law $w$ on $\{\pm1\}^d$. Set
\[
 n=d+1,\qquad N=4\left\lceil\frac n4\right\rceil,
 \qquad C(w)=\sum_xw_x^2.
\]
The integer $N$ is the smallest multiple of four at least $d+1$.
The standard orthogonal-frame argument gives the atom bound
$w_x\le1/n$; the entropy connection and its exact $N=n$ case are
background~\cite{GP}. The residue-sensitive statement below is a
derivation in this research pass; bibliographic priority is unestablished.

\begin{theorem}[Residue-sensitive collision bound and equality]
\label{thm:residue-bound}
For every such law,
\[
 C(w)\le\frac1N,\qquad
 H_2(w):=-\log C(w)\ge\log N,\qquad H(w)\ge\log N.
\]
If $C(w)=1/N$, the law is uniform on exactly $N$ atoms and their
augmented sign matrix admits completion by $N-n\le3$ columns to a
Hadamard matrix of order $N$. Conversely any order-$N$ Hadamard matrix
gives a law attaining both entropy bounds. Consequently attainment
of either entropy lower bound is equivalent to Hadamard existence
at order $N$.
\end{theorem}

\begin{proof}
Write $v_x=(1,x)\in\{\pm1\}^n$. Unbiased pairwise independence says
\begin{equation}\label{eq:residue-frame}
 \sum_xw_xv_xv_x^T=I_n.
\end{equation}
Thus $\sum_yw_y(v_x\cdot v_y)^2=n$ for each $x$, and
$\sum_{x,y}w_xw_y(v_x\cdot v_y)^2=n$.
In particular $w_x\le1/n$.

\emph{Case $n\equiv0\pmod4$.}
The atom bound gives $C\le1/n$. Equality forces every positive weight
to equal $1/n$. The square augmented matrix has orthogonal columns
by \eqref{eq:residue-frame}, and is Hadamard.

\emph{Case $n\equiv3\pmod4$.}
All dot products are odd. Distinct augmented atoms cannot be negatives
of one another, and their squared dot products are at least one.
Consequently
\[
 n\ge n^2C+(1-C),\qquad C\le\frac1{n+1}.
\]
Equality forces every distinct pair to have dot product $\pm1$.
The pointwise identity then reads
$n=n^2w_x+1-w_x$, forcing $w_x=1/(n+1)$.
Let $V$ have these $N=n+1$ augmented rows. Then $V^TV=NI_n$;
$K=NI_N-VV^T$ is positive semidefinite of rank one with diagonal
one. Hence $K=ss^T$ for a sign vector $s$. Appending $s$ gives
$[V\ s][V\ s]^T=NI_N$.

\emph{Case $n\equiv2\pmod4$.}
Here $n\ge6$. Divide the atoms by their parity
$\epsilon_x=\prod_i x_i$; let their total masses be $a,1-a$.
Two atoms in the same class have dot product $2\pmod4$, while
cross-class dot products are $0\pmod4$. Therefore
\[
 n\ge n^2C+4\big(a^2+(1-a)^2-C\big)
       \ge(n^2-4)C+2,
\]
giving $C\le1/(n+2)$. Equality forces $a=1/2$, within-class
distinct dot products of magnitude two, and cross-class orthogonality.
For a fixed atom the pointwise identity is then
$n=n^2w_x+4(1/2-w_x)$, so $w_x=1/(n+2)$.
For the resulting $N\times n$ matrix $V$, the positive semidefinite
residual $K=NI_N-VV^T$ has rank two and diagonal two. Cross-class
blocks vanish, so each of the two nonempty diagonal blocks has rank
one. Write them $2ss^T$ and $2tt^T$, with $s,t$ sign vectors.
Append rows $(s_i,s_i)$ on the first class and $(t_i,-t_i)$ on the
second. These two added columns have row Gram matrix $K$, completing
$V$ to a Hadamard matrix.

\emph{Case $n\equiv1\pmod4$.}
Here $n\ge5$. For any pair of atoms set
$q_{xy}=\epsilon_x\epsilon_y(v_x\cdot v_y)$.
Their Hamming distance shows $q_{xy}\equiv1\pmod4$. Thus
\[
 q_{xy}^2+2q_{xy}-3\ge0,
\]
with equality precisely at $q_{xy}\in\{-3,1\}$.
Put $z=\sum_xw_x\epsilon_xv_x$. Since the $n$ coordinate functions
of $v$ are orthonormal in $L^2(w)$, Bessel inequality gives
$\|z\|^2\le\mathbb E\epsilon^2=1$. Equivalently, this follows by
expanding $\mathbb E(\epsilon-z\cdot v)^2=1-\|z\|^2$.
Summing the polynomial over pairs and retaining its diagonal terms gives
\[
 (n-1)(n+3)C\le n+2\|z\|^2-3\le n-1,
\]
so $C\le1/(n+3)$.

At equality, $\epsilon_x=z\cdot v_x$ on the positive-weight support,
and distinct pairs have $q_{xy}\in\{-3,1\}$. The polynomial's
pointwise expectation equals
$n+2\epsilon_xv_x\cdot z-3=n-1$; its only nonzero contribution
comes from $y=x$. Hence again $w_x=1/(n+3)$.
Switch row $v_x$ to $v'_x=\epsilon_xv_x$ and form $V'$.
Its residual $K=NI_N-V'V'^T$ is positive semidefinite of rank three,
has diagonal three, and off-diagonal entries in $\{3,-1\}$.
Represent it as a Gram matrix of length-$\sqrt3$ vectors. Inner
product three identifies equal vectors. Distinct types have Gram
matrix $4I_t-J_t$, which is positive semidefinite only if $t\le4$.
Assign the types distinct rows of
\[
 T=\begin{pmatrix}1&1&1\\1&-1&-1\\-1&1&-1\\-1&-1&1\end{pmatrix}.
\]
The resulting $N\times3$ sign matrix $R$ has $RR^T=K$.
Then $[V'\ R]$ is a Hadamard matrix. This is an explicit completion;
no external small-rank classification is needed.

For the converse normalize the first column of an order-$N$ Hadamard
matrix to all positive and keep $n$ columns including that column.
The $N$ resulting rows are distinct: full rows differ in $N/2$
coordinates, whereas at most three columns were deleted. For $N\ge8$
this leaves a difference; at $N=4$, $n\ge3$ and at most one column
was deleted. The uniform law on the remaining $d$ coordinates is
unbiased and pairwise independent, with $C=1/N$ and $H=\log N$.
Finally $H\ge-\log C$ by Jensen inequality, so equality in the Shannon
lower bound forces $C=1/N$ as well.
\end{proof}

For example, four, five or six unbiased pairwise-independent sign bits
each have Shannon entropy at least $\log8$. Eight, nine or ten such
bits each have entropy at least $\log12$, attained by truncating an
order-$12$ Hadamard matrix. These statements concern joint entropy,
not min-entropy; a single atom can still saturate $1/(d+1)$.
For $d=1$ the formula would be false, which is why $d\ge2$ is explicit.

\paragraph{The unresolved step.}
This proves a universal bound and a complete equality classification.
It does not prove that an equality law exists in every admissible
dimension. Such attainment for all $d$ is equivalent to the full
Hadamard conjecture. The parity and Bessel argument gives a precise
information-theoretic target; it does not by itself reach that target.
